
The organising question here is \textbf{label quality}. Whether spatial relations can be computed at all is settled. What is not settled is whether computed labels beat the human labels they replace, or how anyone would tell. Rival remedies that stretch scarce labels without replacing them come first (\S{}2.4), and the last two sections argue against the project, on what the field's metrics miss (\S{}2.8) and the case against a rule-based annotator (\S{}2.9).

\textbf{How this chapter was searched.} Google Scholar, ACM DL, IEEE Xplore, OpenReview and arXiv were queried between June and August 2026 on combinations of \textit{scene graph}, \textit{spatial relation}, \textit{automatic annotation}, \textit{pseudo-label}, \textit{robot perception} and the seven predicate names, with forward and backward snowballing from the source paper and from SpatialVLM. Preprints and public code were included, because the systems nearest this one are released that way, so the novelty claim in \S{}1.2.3 is bounded by that search and nothing stronger. Work was kept where it produces or evaluates spatial relation labels. Work where relations stay internal to a model, with no label artefact, was dropped.

\section{Scene graphs and spatial relationships in robotics}

Seven predicates get computed here (\textit{on}, \textit{under}, \textit{left/right of}, \textit{in front of / behind}, \textit{near}), the scene-graph edges carrying the geometry an agent has to respect to act. Told only "cube, book, table", a planner cannot decide what to move first. One told "the cube is on the book" can, and \S{}5.7 turns that into a measurement. Language-driven robot planners ground their instructions in this kind of structured scene state (Ahn et al., 2022). 3D scene graphs were proposed as the unifying form for it (Armeni et al., 2019). Within robotics it has been pushed towards real-time perception, with Kimera building such graphs from SLAM (Rosinol et al., 2021) and 3DSSG from reconstructed indoor scans (Wald et al., 2020). Both take geometry as given and infer the edges. Here the problem runs the other way round, since the geometry has to be recovered from monocular images first, and deciding those edges by rule, which is what I do, is an old idea with a formal literature behind it. Under \textbf{qualitative spatial reasoning}, space is represented through a finite vocabulary of relations, reasoned over with composition tables. RCC-8 fixes connection and containment (Randell, Cui and Cohn, 1992), orientation calculi add direction from a viewpoint (Freksa, 1992), and Cohn and Renz (2008) survey the family. Allen's interval algebra (Allen, 1983) is \textit{temporal}. It only reaches spatial reasoning where its thirteen relations are applied to intervals projected onto an axis.

Three properties of that work matter here. Because the relations are \textit{defined}, what a label means is not in dispute. A definition settles the semantics, not the frame it is read in, nor the perceptual call of whether it holds, which is what \S{}4.5 and \S{}7.3 measure going wrong. Because the calculi are \textit{decidable}, the underdetermined case is where abstention comes from. And the vocabularies are finite and hand-authored, the limitation \S{}2.9 raises and this dissertation concedes. None of it supplies the step from pixels to regions. Section 2.5 covers that, and this project joins the two. So the rest of this review is about how those edges get \textbf{produced}, by hand, by learned prediction, or by computation from measured geometry.

\section{The source dataset and its annotation bottleneck}

Wang et al. (2025) introduce \textit{A Spatial Relationship Aware Dataset for Robotics}, captured by a Boston Dynamics Spot robot. The verified specifics:

\begin{itemize}
\item \textbf{Scale:} "nearly 1,000 robot-acquired indoor images" collected; "approximately 900" survive quality control. (The released GitHub subset used here contains 884 images / \textbf{838 annotated}; \S{}4.1 reconciles the counts.)
\item \textbf{Annotation:} "Nine trained annotators worked independently in batches of 100 images," using the manual \textbf{SGDET-Annotate} tool (every box drawn, every class and relationship clicked), with a majority-vote cleaning pass. The released data is organised into nine groups of \textasciitilde{}100, which matches.
\item \textbf{Vocabulary:} exactly seven predicates (\textit{behind, in front of, on, to the left of, to the right of, under, near}) over six object classes (book, bottle, box, cube, human, remote).
\item \textbf{Detector baseline:} a YOLOv10m backbone (Wang, A. et al., 2024) reaches $\approx$0.92 precision, 0.90 recall and \textbf{0.93 mAP@50} (mAP@50-95 $\approx$0.68).
\end{itemize}

Three limitations the authors flag themselves got me started. Training shows \textbf{early saturation} ("all predictors reached their peak mR@100 well before the final epoch", which they put down to "the dataset's limited diversity"), while \textbf{\texttt{near}} draws "inconsistencies, particularly with the 'near' predicate", which "remains challenging for all models (0.2247--0.2494)". And their future-work prescription is explicit: "augment under-represented relations while enforcing clear annotation guidelines (e.g., spatial thresholds for 'near')." All three come back to the \textbf{manual annotation bottleneck}. SGDET-Annotate only \textit{speeds up} labelling, since a human still decides every edge. So the dataset cannot grow cheaply. Diversity stays low. Disagreement between annotators is baked in. \textbf{No automatic annotator exists for this dataset.} That is the gap this project fills, and the fitted \texttt{near} threshold does directly what the authors asked for.

\section{Label quality: weak supervision and annotator disagreement}

Replacing scarce human labels with dense computed ones already has a name, \textbf{weak supervision}. \textbf{Snorkel} (Ratner et al., 2017) formalised \textit{data programming}, where experts write labelling functions and stop labelling examples one by one, their noisy overlapping votes combining into training labels. The trade is per-label authority for coverage and consistency. Its deployments repeatedly matched or beat hand-labelled baselines wherever the labelled set, not the model, was the bottleneck, which is my position, and my geometric rules are labelling functions in that sense, deterministic, auditable and dense. Two things differ. Measured geometry gives near-exact votes for five predicates (blind-audited 0.79--1.00) and noisy ones for the support pair, where \S{}4.14 puts the shipped rule at 0.54, so the five need no probabilistic combination, and on the other two the problem Snorkel solves by aggregating is solved here by abstaining. I also \textit{check} the computed labels against the human ones they replace (RQ1), so nothing assumes the two are on a level.

Alongside it, a second literature takes apart the idea that human annotation is a single reliable gold standard. Plank (2022) puts it sharply. Variation between annotators is often not error but legitimate difference, and treating it as noise throws away signal and flatters whichever convention the majority held. Chapter 4 reaches that same reading independently for \texttt{near} and front/behind. \textbf{Uma et al.'s (2021) survey of learning from disagreement} documents it across vision and language, driven by ambiguous guidelines, subjective boundaries and annotator-specific conventions. That frame fits what I found. Chapter 4 measures three annotator behaviours, selective \texttt{near} usage, an inverted front/behind convention in two groups, and one-directional support labelling, which between them turn "agreement with the humans" into a per-annotator quantity. Two design decisions follow. I report per annotator group and never only pooled, and thresholds are calibrated only on annotators who used a label, with the rest held out. Where human judgements have to be evaluated, the measurement tradition supplies the instruments. \textbf{Cohen's kappa} (Cohen, 1960) gives chance-corrected agreement between two verdict sets, and Krippendorff's alpha does the same for a pool of raters. Artstein and Poesio (2008) survey both with their pitfalls, including the prevalence effect that depresses kappa when one answer dominates. Chapter 4's independent validation uses these, and sharpens RQ2 into a question the weak-supervision literature predicts but rarely tests: can consistent computed labels \textit{out-teach} inconsistent human ones on the humans' own held-out annotations?

Label quality contaminates \textit{evaluation} too. Northcutt, Athalye and Mueller (2021) measured label errors across ten heavily-used benchmarks (3.3\% average, 6\% of ImageNet's validation labels) and showed that correcting them changes model rankings. From the collection side, spatial-relation benchmarks hit the same problem. SpatialSense (Yang, K., Russakovsky and Deng, 2019) used \textit{adversarial} crowdsourcing, because relations collected without that pressure are dominated by guessable co-occurrences. Rel3D (Goyal et al., 2020) rebuilt the task on 3D scenes with minimally contrastive pairs, having found that 2D datasets let models score well without using spatial information at all. Both respond to something I measure in this dataset too. What a benchmark appears to test and what its annotation rewards can drift apart, and nobody notices until someone checks the labels. Chapter 6 tests that, and the consequence takes a different shape from Northcutt et al.'s, since correcting the gold lifts both label sources together without reordering them, and what moves with annotation quality is which source leads on which annotator.

\section{The rival family: semi-supervised and active learning}

Computing labels from geometry is not the only answer to expensive annotation, and the stronger rival, which I run as an arm of the experiment, is to take the labels that already exist and stretch them. \textbf{Active learning} (Settles, 2009) cuts annotation cost by choosing \textit{which} examples a human labels next. \textbf{Semi-supervised learning} trains on a small labelled set plus a large unlabelled one (van Engelen and Hoos, 2020). Its commonest instrument is \textbf{pseudo-labelling} (Lee, 2013), where a model trained on the labelled seed labels the rest for its own retraining, its strongest modern form being noisy self-training (Xie et al., 2020). Applied here, the recipe would be to train on the \textasciitilde{}10\% of pairs the annotators labelled, pseudo-label the remaining 90\%, then retrain. Three measured properties of this dataset argue against it.

Self-training \textit{amplifies its seed}, and this seed is sparse and internally inconsistent (selective \texttt{near}, two inverted front/behind conventions, one-directional support; \S{}2.2, Chapter 4). The seed is also \textit{selectively} small, since annotators labelled what they found salient, so the labelled 10\% is not the unbiased sample of the 90\% pseudo-labelling assumes. And Chapter 5's human-trained classifier is exactly the seed such a loop would start from. It collapses on the sparsely-labelled predicates, where recall runs 0.08--0.25, and is unstable across seeds, which leaves little reliable enough to be worth amplifying. Active learning fails differently. It still buys \textit{human} labels, so it makes the bottleneck shallower without removing it, and it rations inconsistency without fixing it. The geometric route avoids all three, because its labelling function does not come from the flawed seed. The rules are fitted to a handful of thresholds, checked on held-out annotators, and are as consistent on the 90\% as on the 10\%. None of this is left as argument. Chapter 5 runs the rival as a third arm of the controlled experiment, with the standard teacher-student loop over the same features, model, split and seeds, so pseudo-labelling and programmatic labelling meet head to head on the humans' own held-out annotations. Both literatures motivate that comparison and neither usually runs it.

\section{Geometry-to-label pipelines (the work this builds on)}

In one line of work the spatial facts are \textit{computed} from perceived geometry, with nothing predicted from learned patterns, and the idea goes back further than most citations suggest. CLEVR (Johnson et al., 2017) emitted the spatial relations of hundreds of thousands of scenes \textit{from the renderer}, exact by construction, and became the standard diagnostic for compositional reasoning because programmatic labels carry no annotator noise to memorise. It also sidesteps the hard half, since its geometry is known because the scenes are synthetic, while an annotator for real photographs has to recover geometry from pixels first. Each of the works below takes up that half. Each outputs something other than a scene-graph annotation for these seven predicates, and the assessments are my own, against the requirements in \S{}2.10.

\begin{itemize}
\item \textbf{SpatialVLM} (Chen et al., 2024), the foundational geometry-to-label method, lifts internet images to metric 3D via monocular depth and segmentation and emits up to \textasciitilde{}2B spatial question-answer pairs from \textasciitilde{}10M images. It establishes the premise this project adopts: spatial relations can be \textit{derived from measured geometry without human relational labels}. Its output is \textbf{free-form VQA text}, against no fixed predicate set.
\item \textbf{SpatialRGPT} (Cheng et al., 2024) adds a curation pipeline that learns regional representations from 3D scene graphs, plus a \textbf{depth plugin} for a VLM's visual encoder. It is the cleanest published RGB$\rightarrow$depth$\rightarrow$relation recipe and the key reference for this project's depth use, but its artefact is a region-reasoning VLM, not a labeller writing VG-format triplets.
\item \textbf{VQASynth} (Remyx AI, 2024) reproduces that pipeline openly (SAM2, monocular depth, grounded captioning). It is the most reusable code reference, but it produces QA pairs where a triplet writer is needed, and its depth backend has shifted, so this project pins its own.
\item \textbf{Open3D-VQA} (Zhang et al., 2025) is where the \textbf{error-correction} idea comes from: extraction of 3D spatial relationships from a single RGB image with a correction flow that discards whatever its own geometry declares impossible. Section 3.6 adapts it for a different domain and output.
\item \textbf{RoboSpatial} (Song et al., 2025) is the closest robotics-domain match and is cited by the source paper. It teaches spatial understanding to 2D/3D VLMs from real indoor scans and formalises \textbf{reference frames}: ego-centric, world-centric and object-centric readings of one phrase. That ambiguity is a documented property of spatial language, not an engineering nuisance. Landau and Jackendoff (1993) showed language encodes location through frame-dependent primitives, so "the cup is left of the box" is true in one frame and false in another, and an annotator has to pin the frame before any label is well defined. Hence \textit{left/right} in the \textbf{camera frame} (Supplementary C), which is the frame the dataset's annotators saw on screen. Its output is spatial QA over three frames.
\end{itemize}

The obvious alternative is to ask a capable vision-language model to name the relations, and there is reason to doubt it before testing. Visual Spatial Reasoning finds a wide model-human gap on relations a person reads off instantly (Liu, Emerson and Collier, 2023), and Kamath, Hessel and Chang (2023) show the failure survives scale and prompting. Section 4.13 tests it. Both models recover under half the human triplets the pipeline does and lose F1 on every predicate, while being \textit{more precise} where they do speak. What settles it is the shape of the output: silence on most pairs, and a symmetric relation asserted in one direction only about a third of the time. Those are the behaviours \S{}4.5 measures in the \textit{human} annotation. Asked to annotate, a vision-language model reproduces the failure mode this project set out to replace. One adjacent family needs separating out, because from robotics it looks closest. Online 3D scene-graph \textit{mapping} systems build a spatial-semantic graph as a robot moves, with Hydra doing it from depth-equipped SLAM in real time (Hughes, Chang and Carlone, 2022) and ConceptGraphs by fusing foundation-model features into an RGB-D map (Gu et al., 2024). They consume depth sensors, emit no dataset-format annotation for existing monocular images, and are not validated against human annotators. If anything they strengthen the case for automatic annotation, since their learned components need the training data an annotator supplies.

Across this body of work, almost nobody asks how anyone knows the computed labels are right, which matters here because my central claim is a validation claim. Two kinds of evidence get offered, and neither is what RQ1 needs. The first is \textbf{downstream benefit}, where SpatialVLM fine-tunes a model on its supervision and shows better answers, and SpatialRGPT judges its representations through the model they produce. The second is \textbf{internal consistency}, where Open3D-VQA discards whatever its own rules declare impossible. Both share a blind spot. A model trained on computed labels and tested on questions drawn from the same computation will score well on any convention applied consistently, including a wrong one, and an internal consistency check is satisfied by any convention that hangs together. Neither detects systematic disagreement with how humans use the words, the failure mode Chapter 4 measures in the tool and the annotators alike. Only comparison against an independently produced human record can, none of these works does it, and that comparison is RQ1. It also explains why Chapter 6's less favourable result does not contradict Chapter 5's: judged against \textit{human} annotation, downstream benefit becomes genuinely adversarial, and this body of work does not run that test.

\textbf{Synthesis.} Every pipeline here targets a \textit{different output} (VQA text or a live map), a \textit{different domain or sensor suite} (internet images, scans, RGB-D), or is a \textit{reference recipe}, not a deployable annotator. As an idea, geometry-to-label is well established. Its application as a \textbf{fully-automatic seven-predicate annotator validated against this dataset's human labels} is not.

\section{Perception components and the open-vocabulary family}

Although the pipeline is geometry-first, it stands on off-the-shelf perception, and I chose each component against alternatives. This design space is also the basis of the \textbf{detector-swap ablation}:

\begin{itemize}
\item \textbf{Detection.} \textit{Closed-set:} \textbf{YOLOv10m} (Wang, A. et al., 2024), the dataset's own baseline at 0.93 mAP@50, is what a replication with the authors' weights would use. \textit{Open-vocabulary:} \textbf{Grounding DINO} (Liu et al., 2024) trades per-class sharpness for arbitrary text-named objects, which carries the pipeline beyond the six classes. Chapter 4's deployment mode uses it at a 0.25 box threshold precisely because it is the worst reasonable detector, so the end-to-end bound it produces (0.38 triplet recall, against 0.85 conditional on both endpoints being found) errs low. Both extend the transformer-detector line DETR began by removing hand-designed anchors and non-maximum suppression (Carion et al., 2020).
\item \textbf{Segmentation.} \textbf{SAM2} (Ravi et al., 2024), box-prompted, the video successor to Segment Anything. It was the earlier model whose promptable formulation and eleven-million-image corpus made class-agnostic segmentation a component other systems simply call (Kirillov et al., 2023); SAM2 inherits that formulation and adds video. The silhouette does work twice over: depth is sampled by median over object pixels, and the support rule's contact test needs the object's bottom boundary pixel by pixel (\S{}3.5).
\item \textbf{Monocular depth.} \textbf{Depth Anything v2} (Yang, L. et al., 2024) emits a \textit{relative} map, ordering pixels without a unit, and that property fixes the rule design: every depth comparison is ordinal and within-image. The scale ambiguity comes with the zero-shot setting. MiDaS established the mixed-dataset training that makes cross-domain monocular depth work, and paid for it with predictions defined only up to an unknown scale and shift (Ranftl et al., 2022).
\item \textbf{Adjacent, not used.} \textbf{CLIP} (Radford et al., 2021) and \textbf{SCLIP} (Wang, F. et al., 2024) bear on open-vocabulary scaling rather than relation logic; \textbf{PrimitiveAnything} (Ye et al., 2025) assumes clean 3D input this project does not have. All three are future directions.
\end{itemize}

The split of work is deliberate. The neural components only \textit{measure}, and every relationship decision is an explicit rule over those measurements, which makes the annotator auditable. Section 3.4 gives each choice with the alternative it displaced, and Chapter 4's box-only ablation measures what each one contributes.

\section{Learned scene-graph generation (the consumer of our output)}

Scene-graph generation (SGG) models \textbf{predict} relationships from learned visual patterns, and the task is older than the framing (Lu et al., 2016, whose language priors already leaned on label statistics). The first whole-graph models passed messages between object and relationship nodes until the two agreed (Xu et al., 2017), and the field's shape was set by \textbf{Visual Genome} (Krishna et al., 2017), 108k images of crowdsourced triplets whose JSON format this dataset inherits. Chang et al. (2023) survey the background and name annotation cost and label bias as its two persistent constraints, which is this project's premise seen from the consumer's side.

\textbf{Neural Motifs} (Zellers et al., 2018) showed that context and label statistics dominate. Their frequency baseline, predicting the commonest predicate while ignoring the image, proved hard to beat, a warning that relation "accuracy" can be measuring memorised co-occurrence. \textbf{VCTree} (Tang et al., 2019) composes dynamic tree structures and is the best model in the source paper's own benchmark (mR@100 = 0.49). \textbf{Unbiased SGG} (Tang et al., 2020) showed formally that such models absorb the \textit{annotation distribution}, and proposed counterfactual debiasing. Where the field went next is telling, since panoptic scene graph generation (Yang, J. et al., 2022) replaced box grounding with pixel-accurate masks, having shown that boxes systematically mislocalise the objects whose relations are being learned, the same reasoning that puts SAM2 masks at the centre of this project's support rule.

This work also fixes the evaluation vocabulary used in Chapters 4 and 7. SGG models are scored by recall of annotated triplets among their top K predictions. \textbf{R@K} pools predicates, so frequent ones dominate. \textbf{mR@K} averages per predicate. \textbf{zR@K} scores only triplet types never seen in training (Tang et al., 2020). Two settings matter. \textbf{PredCls} supplies ground-truth boxes and classes, which isolates the relation model, while \textbf{SGDet} requires detection, labelling and relations end to end, and I adopt both conventions so that my numbers read against the source paper's tables. Two lessons carry over. The models consuming this dataset's labels are known bias-absorbers, so whatever the annotation carries becomes the training signal. Cleaning the supply goes after the cause; debiasing the model only treats the symptom. And the frequency-baseline lesson from Motifs sets the baseline discipline: every fidelity number in Chapter 4 is read against trivial random and majority baselines.

\textbf{REACT++} (Neau and Falomir, 2026), a preprint at the time of writing, is a real-time SGG model with a YOLO backbone, reportedly \textasciitilde{}20\% faster and \textasciitilde{}10\% more accurate on relation prediction than its predecessor, small enough to run onboard a robot, and shipped in the open \textbf{SGG-Benchmark} framework. What matters is that models like this \textbf{need labelled training data and sit downstream of an annotator}. This project computes labels and runs \textit{before} any SGG model, which makes it the \textbf{supplier} and REACT++ a natural \textbf{consumer}. Training REACT++ on the auto-labels versus the human labels is the heavyweight version of RQ2, run in Chapter 6.

\section{How the field measures success, and what those measures miss}

Every metric in \S{}2.7 was adopted to fix a defect in the one before it, and each carries a defect of its own. Chapters 4 and 6 report results \textit{in} these metrics and then argue about what they mean, so the arguments belong here, grounded in the literature.

\textbf{Recall gets used without precision because the annotation is incomplete.} The convention comes from visual relationship detection on crowdsourced graphs (Lu et al., 2016; Krishna et al., 2017), where annotators record a handful of the relations present. An unannotated pair has not been examined, which is not the same as being negative, so a predicted relation absent from the gold cannot be scored wrong, and the field dropped precision and ranked by recall at K. What it costs is a metric that cannot tell careful prediction from abundant prediction, and that bites hardest when the annotation is itself the object of study. A method labelling the pairs humans skipped is penalised for coverage under precision and rewarded under recall, and neither settles whether the extra labels are true, which neither constraint Chang et al. (2023) name (\S{}2.7) resolves.

\textbf{A metric that a context-free prior can saturate is measuring the prior.} Zellers et al.'s (2018) frequency baseline proved hard to beat on R@K, which says as much about the metric as about the models. \textbf{mR@K}, proposed by VCTree (Tang et al., 2019) and standard after Tang et al. (2020), is the correction, and brings its own sensitivity, since a predicate with few test instances weighs as much as one with thousands, so the total can move on a handful of triplets, and in a dataset annotated in blocks by different people, on which annotator supplied them. Per-predicate and per-annotator breakdown is therefore not decoration. Without it a mean recall cannot be interpreted, which is why Chapters 4 and 6 report both.

\textbf{Zero-shot recall measures something relative, and what it is relative to has to be stated.} zR@K scores only combinations absent from \textit{training}. It comes from visual relationship detection (Lu et al., 2016) and was first reported on Visual Genome by Tang et al. (2020). It is well posed when a model is compared against itself. The moment two differently-supervised models share a column it turns ambiguous, because the exclusion set then comes from one shared reference, and what it reports is how far each label source covers combinations the reference omits. That is a real property, arguably the more interesting one for an annotation study, but it is not compositional generalisation. Section 6.4 shows the distinction is not hypothetical here, and reports the quantity under its accurate name.

\textbf{Setting and gold determine what a number means.} PredCls produces systematically higher figures than SGDet, where detection errors propagate. Quoted side by side, as they often are, the two are not comparable. Underneath both sits the assumption that the gold is correct, which Northcutt, Athalye and Mueller (2021) showed is false often enough to reorder published rankings (\S{}2.3). For spatial relations the annotation may also be \textit{consistently} wrong in the way that matters here. Where annotators applied a different reference frame, a system agreeing with them scores well and a correct system scores badly, and no recall metric can tell the two apart.

Three commitments follow. I report recall against the human triplets alongside an audited estimate of precision on pairs the gold never covers (\S{}4.4). Every total is broken down per predicate and per annotator group. And the decisive test of the labels is deliberately moved \textit{off} these metrics altogether, onto whether a downstream consumer trained on them does the task (Chapters 5 and 6), because a metric that rewards agreement with the annotation cannot settle a dispute about the annotation.

\section{The case against a rule-based annotator}

Below are the strongest objections to the method, not the easiest ones to answer, four against the method and a fifth against its premise, set out before any result so the later chapters can be read as attempts on them.

\textbf{Rules do not scale with the vocabulary.} Each predicate is an explicitly written geometric test with fitted thresholds. Seven are manageable. The literature routinely works with fifty (Krishna et al., 2017), and the direction of travel Chang et al. (2023) identify is open-vocabulary relations, including non-spatial ones such as \textit{holding} for which no geometric criterion exists. A learned predictor improves by being shown more data, whereas a rule set only improves when somebody extends it by hand, a difference in kind, not in degree. So whatever this project shows about seven spatial predicates transfers to functional relations not at all.

That boundary, which took me a while to place, is worth locating precisely, because there is a literature on the far side of it. \textit{Holding} is a \textbf{human-object interaction}, and detecting those is a task in its own right with its own benchmarks. Chao et al. (2018) frame it as predicting an interaction label over a human-object pair, learned from examples, with geometry deciding nothing. Underneath sits the older notion of an \textbf{affordance}, what an object offers an agent, introduced by Gibson (1979) and surveyed by Jamone et al. (2018) across psychology, neuroscience and robotics. The recurring finding is that what an object affords depends on the agent and its intent, so it cannot be recovered from shape and position alone, and this project meets the same wall from the other side. Section 3.6 cannot separate a person \textit{holding} an object from a surface \textit{supporting} one by geometry, and ablation A10 (D.8) measures that failure, so the pipeline excludes the \texttt{human} class from support by configuration. The literature says a case like this gets decided by learned interaction evidence, not by a better geometric test, which is why the exclusion marks a boundary of the method, not a defect in how I built it.

\textbf{Systematic error is worse for training than random error.} Consistency only guarantees that the mistakes recur, and a rule's mistakes, unlike a person's, correlate with scene geometry where independent noise would scatter. Tang et al. (2020) established how thoroughly SGG models absorb the distribution of their supervision, and that result cuts both ways: a model trained on rule output can learn the rule's blind spot as a fact about the world, and no amount of extra data averages it out. In this narrow respect, independent human noise is the safer failure mode.

\textbf{Validating one's own labelling functions is circular.} Snorkel (Ratner et al., 2017) anticipated this, treating labelling functions as noisy and \textit{estimating} their accuracies from the agreement structure among several independent sources, precisely because an author's confidence in a rule is not evidence about it. A single rule set verdicted by its own author has neither, and an audit by the person who wrote the rules inherits their assumptions about what counts as correct. Chapter 4 concedes it. The mitigations are reported there with their limits.

\textbf{The reference frame is a decision, not a fact.} \textit{In front of} has no answer independent of the frame it is asked in. Landau and Jackendoff (1993) distinguish viewer-, object- and environment-centred description, and RoboSpatial (Song et al., 2025) maintains all three rather than choosing. A rule set has to commit to a convention, and where an annotator used a different one the two will disagree systematically. Calling the rule correct there is an assertion about which convention should govern, not a measurement, and the dissertation has to argue for it rather than assume it (\S{}4.5).

Aimed at the premise, the fifth objection is different in kind. If the existing annotation is inconsistent, the direct fix is better collection, not cheaper labels. SpatialSense (Yang, K., Russakovsky and Deng, 2019) and Rel3D (Goyal et al., 2020) both responded to defective relation annotation by rebuilding the collection process. Automation makes labels cheap. Getting the definitions right is a separate job, and I inherited this dataset's definitions without improving them. The dissertation answers the second and fourth objections empirically (Chapters 5--6, and Chapters 4 and 7 respectively), concedes the third while reporting what mitigation was possible, does not answer the first at all and records it as a limitation (\S{}8.3). Answering the fifth would be a different project.

\section{Critical comparison and the research gap}

Every neighbour below targets a different output, is a reference recipe, not an annotator, works outside this dataset, or \textit{predicts} instead of \textit{computing} relations. \textbf{None provides a fully-automatic annotator for this dataset's seven predicates, validated against its human labels.}

\needspace{16.8\baselineskip}
\begin{center}
\setlength{\tabcolsep}{3pt}
\scriptsize
\begin{longtable}{>{\raggedright\arraybackslash}p{0.126\textwidth}>{\raggedright\arraybackslash}p{0.094\textwidth}>{\raggedright\arraybackslash}p{0.155\textwidth}>{\raggedright\arraybackslash}p{0.156\textwidth}>{\raggedright\arraybackslash}p{0.153\textwidth}>{\raggedright\arraybackslash}p{0.110\textwidth}>{\raggedright\arraybackslash}p{0.103\textwidth}}
\caption[Related work assessed against the requirements]{Related work assessed against the requirements of a fully-automatic spatial-relationship annotator.}\\
\hline
Work & Year/venue & Relations: compute vs. predict & Output & Auto-annotator? & This dataset's 7 predicates? & Validated vs. these human labels? \\
\hline
\endfirsthead
\hline
Work & Year/venue & Relations: compute vs. predict & Output & Auto-annotator? & This dataset's 7 predicates? & Validated vs. these human labels? \\
\hline
\endhead
Wang et al. (source) & 2025 ACM MM & human-labelled & VG JSON / YOLO / h5 triplets & manual (SGDET-Annotate) & defines them & is the human label \\
SpatialVLM & 2024 CVPR & compute (geometry) & spatial VQA text & partial (QA gen) & no & no \\
VQASynth & 2024 & compute (geometry) & spatial VQA text & partial (QA gen) & no & no \\
SpatialRGPT & 2024 NeurIPS & compute (geometry+depth) & region-reasoning VLM & reference recipe & no & no \\
Open3D-VQA & 2025 ACM MM & compute (+correction) & 3D/aerial VQA & partial & no & no \\
RoboSpatial & 2025 CVPR & compute (geometry) & spatial QA (3 frames) & partial & no & no \\
REACT++ / SGG-Benchmark & 2026 & \textbf{predict} (learned) & scene-graph triplets & no (needs labels) & no & n/a \\
\textbf{This work} & 2026 & \textbf{compute (geometry)} & \textbf{VG JSON / YOLO / h5 triplets} & \textbf{yes, fully automatic} & \textbf{yes} & \textbf{yes (RQ1)} \\
\hline
\end{longtable}
\end{center}

Two columns isolate the contribution. Only Wang et al. and this work address the seven predicates, and Wang et al. do it manually. Only this work quantifies agreement with the human consensus on the same images. The geometry-to-label \textit{method} is borrowed and well-precedented; building it into a validated automatic annotator for this dataset is new.

\section{Summary and positioning}

The gap, once all of that is set against each other, is a narrow one. No automatic, geometry-based annotator emits these seven predicates in the dataset's native formats and is validated against its human labels, even though its authors ask for exactly that. Two constraints follow. The field's metrics are recall-shaped by its own incomplete annotation (\S{}2.8), so no number from them can settle a dispute \textit{about} that annotation, and the protocol of Chapters 4 to 6 is built around that limit, not inside it. And the strongest case against the approach is stated in advance (\S{}2.9), so the results can be read as an attempt on it, with \S{}7.7 the reckoning. Chapter 3 turns this gap into a design: the methodology that structures the work, and the geometric rule for each of the seven predicates.